\documentclass[12pt, letterpaper]{article}
\usepackage[margin=1in]{geometry}
\usepackage{amsmath}
\usepackage{graphicx}
\usepackage{booktabs}
\usepackage{hyperref}
\usepackage{float}
\usepackage{caption}
\usepackage{xcolor}
\usepackage{array}
\usepackage{subcaption}

\title{\textbf{In-Situ Balling Detection in Laser Powder Bed Fusion\\
Using Multi-Modal Camera and Surface Profilometry}}
\author{Erfan Ziad}
\date{\today}

\begin{document}
\maketitle

\begin{abstract}
Balling is a common defect in Laser Powder Bed Fusion (LPBF) in which
the molten track breaks up into discrete spherical beads rather than
forming a continuous weld. We present a multi-modal machine learning
approach to detect balling in-situ using a top-down melt pool camera
combined with post-build surface profilometry. A CNN-LSTM model
processes sequences of consecutive melt pool image differences and,
for layers with available profilometry, a spatially aligned 2D height
patch extracted from the surface height map. Model performance is
evaluated using Leave-One-Layer-Out (LOLO) cross-validation across
14 single-track experiments spanning two scan speeds and six power
levels. The best model achieves mean LOLO AUC of 0.720 compared to
0.673 for a camera-only baseline, with a mean improvement of 0.126
AUC on the five layers where surface height data is available.
Key findings include: bead position type is the primary determinant
of camera detectability, driven by a fixed gas flow asymmetry;
camera and profilometry capture complementary physical phenomena
with near-zero cross-correlation ($r = -0.056$); and Middle-type
beads and small beads previously undetectable from camera alone
become detectable when height data is incorporated.
\end{abstract}

% ══════════════════════════════════════════════════════════════════
\section{Introduction}
% ══════════════════════════════════════════════════════════════════

Laser Powder Bed Fusion (LPBF) is a metal additive manufacturing
process in which a laser selectively melts successive layers of
metal powder to build complex three-dimensional parts. Despite its
versatility, LPBF is susceptible to process instabilities that
produce defects affecting mechanical properties and dimensional
accuracy. Balling is one such instability: when surface tension
forces overcome the laser-driven wetting of the melt track, the
molten pool breaks into discrete spherical beads rather than
forming a continuous solid track. Left undetected, balling causes
surface roughness, inter-layer bonding failures, and in severe
cases, build termination.

In-situ detection of balling is challenging because the defect
manifests at the scale of the melt pool ($< 1$\,mm) at high scan
speeds ($\sim 2000$\,mm/s) with no ground truth available during
the build. Most deployed sensing approaches use a single top-down
camera to capture the melt pool emission. While such cameras
provide real-time spatial and temporal information, they are
fundamentally limited to what is optically detectable — pool
dynamics and emission intensity — which does not always correspond
to the post-build bead geometry measurable by profilometry.

This paper makes the following contributions:
\begin{enumerate}
\item A systematic characterization of which balling categories
are detectable from a top-down melt pool camera and why, including
a physical explanation of the observed Left/Right detectability
asymmetry through gas flow geometry.
\item A multi-modal CNN-LSTM architecture that fuses melt pool
image sequences with spatially aligned profilometry height patches,
with an explicit gating mechanism for layers where height data
is unavailable.
\item An ablation study demonstrating that camera and profilometry
signals are genuinely complementary ($r \approx 0$) and that
their fusion substantially improves detection of previously
undetectable bead categories.
\item A process parameter ablation showing that scan speed $V$
conditioning can hurt generalization through shortcut learning,
while laser power $P$ provides a more useful continuous signal.
\end{enumerate}

% ══════════════════════════════════════════════════════════════════
\section{Dataset}
% ══════════════════════════════════════════════════════════════════

The dataset consists of 14 single-track LPBF experiments conducted
at two scan speeds (V = 1800 and 2000\,mm/s) across six power
levels (P = 320--460\,W). Each layer was imaged by a top-down
melt pool camera at high frame rate, and after building, each
layer was inspected by optical microscopy to label each camera
frame as balling or clean. Balling events are further annotated
by bead position type (Left, Middle, Right relative to the track
centerline) and bead size (Large $\geq 0.15$\,mm, Medium
0.10--0.15\,mm, Small $< 0.10$\,mm).

\begin{table}[H]
\centering
\caption{Dataset summary. V = 1800\,mm/s produces consistently
lower balling rates (1.6--28.6 percentage points) across all
matched power levels.}
\begin{tabular}{lcccc}
\toprule
Group & V (mm/s) & P range (W) & Total frames & Balling frames \\
\midrule
L226--L231 & 1800 & 360--460 & 1{,}374 & 260 (18.9\%) \\
L245--L252 & 2000 & 320--460 & 1{,}704 & 510 (29.9\%) \\
\midrule
Total & --- & 320--460 & 3{,}078 & 770 (25.0\%) \\
\bottomrule
\end{tabular}
\end{table}

Five layers (L248--L252, V = 2000\,mm/s) additionally have
post-build surface profilometry: a $500 \times 1170$ pixel
height map at anisotropic resolution ($p_x = 7.89\,\mu$m,
$p_y = 12.50\,\mu$m).

% ══════════════════════════════════════════════════════════════════
\section{Camera Signal and Detectability}
% ══════════════════════════════════════════════════════════════════

\subsection{Raw Signal Strength}

Before model training we measured Cohen's $d$ — the normalized
difference in mean frame-to-frame pixel change between balling
and clean frames — as a proxy for raw camera signal strength
per layer (Figure~\ref{fig:cohens_d}).

\begin{figure}[H]
\centering
\includegraphics[width=\linewidth]{a2_signal_cohens_d.png}
\caption{Cohen's $d$ per layer. Positive values indicate balling
frames are more visually active than clean frames.
V = 1800\,mm/s layers (red) show stronger and more consistent
signal; V = 2000\,mm/s layers (blue) show weak or negative signal
in most cases.}
\label{fig:cohens_d}
\end{figure}

V = 1800\,mm/s layers show substantially stronger signal, with
L227 (P = 380\,W) reaching $d = +0.646$ — approximately
$3\times$ stronger than any V = 2000\,mm/s layer. Most
V = 2000\,mm/s layers have near-zero or negative Cohen's $d$,
indicating that balling at higher scan speeds produces subtler
camera perturbations, consistent with shorter melt pool residence
time and less pronounced pool dynamics.

\subsection{Detectability by Bead Type and Size}

From combined LOLO predictions of the camera-only baseline,
bead \textbf{position type} emerges as the primary determinant
of detectability, not bead size.

\begin{table}[H]
\centering
\caption{AUC by bead type and size from camera-only baseline.
Pooled across all 14 LOLO folds. AUC $> 0.65$ considered
detectable.}
\begin{tabular}{llccc}
\toprule
Type & Size & N & AUC & Detectable? \\
\midrule
Right  & Large  & 21  & 0.853 & Yes \\
Right  & Medium & 45  & 0.735 & Yes \\
Middle & Large  & 15  & 0.793 & Yes \\
Middle & Medium & 109 & 0.652 & Yes \\
Left   & Large  & 23  & 0.744 & Yes \\
Left   & Medium & 51  & 0.674 & Yes \\
Left   & Small  & 84  & 0.756 & Yes \\
Right  & Small  & 27  & 0.509 & No \\
Middle & Small  & 381 & 0.640 & No \\
\bottomrule
\end{tabular}
\end{table}

\textbf{Physical interpretation.} Lateral beads (Left and Right)
displace the melt pool asymmetrically, creating a transient
bright structure in consecutive frame differences that the
camera reliably detects. Middle beads form directly under
the laser path and are remelted symmetrically — no lateral
signal appears in difference images, a fundamental limitation
of the single top-down camera configuration. Small beads
($< 0.10$\,mm) cause sub-resolution perturbations
indistinguishable from normal pool dynamics.

\subsection{Gas Flow Asymmetry Explains Left vs Right Detectability}

Right-type beads (AUC = 0.804 overall) are consistently more
detectable than Left-type beads (AUC = 0.643). This asymmetry
has a direct physical explanation rooted in the fixed process
geometry: the laser scans $-x \to +x$ while the shielding gas
flows $+y \to -y$, perpendicular to the scan. The right side
of the track is persistently downwind — shielding gas carrying
spatter, metal vapor, and ejecta arrives at the melt pool from
the right. Right-type beads therefore form in a more disturbed
local atmosphere, producing stronger and more consistent pool
perturbations visible to the camera. Left-type beads form
upwind in cleaner gas conditions where the bead-pool
interaction is subtler and more variable across layers.

This is a genuine physical finding rather than a camera artifact:
the CNN-LSTM model detects a real asymmetry in melt pool dynamics
driven by fixed scan and gas flow geometry. For future experiments,
alternating scan direction between layers would eliminate this
systematic bias and allow unbiased comparison of Left and Right
bead detectability.

\subsection{Left-Type Bead Inconsistency}

Left-M bead AUC varies widely across layers (0.398--0.816).
Pixel-level analysis of the horizontal centroid displacement of
the brightest difference signal reveals Pearson $r = -0.545$
between centroid leftward displacement and per-layer AUC
(Figure~\ref{fig:leftm}): more clearly lateral beads are more
detectable. The ``Left'' microscopy label encompasses beads
with a wide range of actual lateral displacements — some
near-center beads appear visually identical to Middle beads
in the camera. This label impurity, rather than training data
quantity ($r = -0.295$) or camera geometry, is the primary
explanation for Left-M inconsistency.

\begin{figure}[H]
\centering
\includegraphics[width=\linewidth]{leftM_summary.png}
\caption{Left-M detectability analysis. Centroid displacement
from track center correlates with per-layer AUC ($r = -0.545$):
more laterally displaced beads are more detectable.
Green bubbles (detectable) cluster on the left; red near center.}
\label{fig:leftm}
\end{figure}

% ══════════════════════════════════════════════════════════════════
\section{Method}
% ══════════════════════════════════════════════════════════════════

\subsection{Camera Representation}

We represent each camera observation as a sequence of consecutive
frame differences $d_t = |I_t - I_{t-1}|$ over a sliding window
of SEQ\_LEN = 8 frames. Differencing isolates transient pool
perturbations (balling events) from the static background,
and the sequence provides temporal context for events that unfold
over multiple frames. Balling events in this dataset have a
median duration of 6 frames and a 75th percentile of 7 frames;
a window of 8 frames (spanning $\approx 0.32$\,mm at V = 2000\,mm/s)
covers the majority of events with context on each side.

An input representation ablation (Table~\ref{tab:input_ablation})
comparing raw frames, difference images, and their concatenation
confirmed that difference images better detect lateral bead types
(Left AUC gap: $+0.174$; Right: $+0.131$) while raw frames are
marginally better for Middle-type beads. Since lateral bead
detection is the primary practical goal and DIFF outperforms
RAW on all bead type strata in pooled evaluation, difference
images are used throughout.

\begin{table}[H]
\centering
\caption{Input representation ablation — per-stratum pooled AUC.
DIFF = consecutive differences. RAW = raw frames. CONCAT = both.}
\label{tab:input_ablation}
\begin{tabular}{lcccc}
\toprule
Stratum & DIFF & RAW & CONCAT & Best \\
\midrule
Overall     & \textbf{0.645} & 0.535 & 0.591 & DIFF \\
Type=Middle & 0.573 & 0.492 & \textbf{0.574} & CONCAT \\
Type=Left   & \textbf{0.770} & 0.596 & 0.628 & DIFF \\
Type=Right  & \textbf{0.877} & 0.746 & 0.751 & DIFF \\
Size=S      & \textbf{0.605} & 0.507 & 0.556 & DIFF \\
\midrule
Mean AUC    & 0.648 & \textbf{0.673} & 0.653 & \\
\bottomrule
\end{tabular}
\end{table}

\subsection{Surface Profilometry Alignment}

Post-build surface profilometry is available for five layers
(L248--L252). Aligning the height map coordinate system with
the camera frame coordinates requires two offsets: $X_\text{off}$
(along scan) and $Y_\text{off}$ (cross-track). $Y_\text{off}$
is calibrated from the track pitch — five bead ridges in the
height map are separated by the known 1\,mm machine pitch.
$X_\text{off}$ is refined per layer via cross-correlation
between the binary balling label signal (converted to a
position-indexed binary vector) and the profilometry height
profile extracted from a $\pm 0.25$\,mm band around each
track centerline.

\begin{figure}[H]
\centering
\includegraphics[width=\linewidth]{full_overlay.png}
\caption{Surface height map with aligned laser scan paths.
Each colored line sits on its corresponding bead ridge,
confirming spatial alignment between camera coordinates and
profilometry coordinates.}
\label{fig:overlay}
\end{figure}

\begin{figure}[H]
\centering
\includegraphics[width=\linewidth]{along_scan_label_spans_aligned.png}
\caption{Along-scan height profiles with balling label spans
overlaid. Colored bands indicate balling events by type;
red dots mark individual balling frame positions.
Labels show bead type, size category, and diameter in mm.}
\label{fig:spans}
\end{figure}

\subsection{Camera-Profilometry Independence}

A critical prerequisite for multi-modal fusion is that the two
modalities provide non-redundant information. The Pearson
correlation between local max profilometry height at each frame
position and the camera-only model's prediction probability is
$r = -0.056$ ($p = 0.076$, $N = 995$) — not statistically
significant and close to zero.

Physically, this independence is expected: the camera detects
real-time pool dynamics during laser scanning (lateral asymmetry,
spatter, emission intensity), while the profilometry measures
the final solidified surface geometry after the scan is complete.
These are different physical quantities measured at different
times. Strikingly, Right-type beads — the most camera-detectable
category — have \emph{lower} profilometry height than clean
frames (Cohen's $d = -0.378$): the camera detects pool
perturbation driven by gas flow, not bead height. Middle-type
beads are slightly taller than clean on average ($d = +0.193$)
but the effect is small and inconsistent across layers.

\begin{figure}[H]
\centering
\includegraphics[width=\linewidth]{height_distributions.png}
\caption{Profilometry height distributions by bead type and size.
Distributions substantially overlap between balling and clean
frames, consistent with the near-zero correlation with camera
prediction scores.}
\label{fig:height_dist}
\end{figure}

\subsection{Model Architecture}

The model (Figure~\ref{fig:arch}) processes four inputs
simultaneously for each 8-frame camera window:

\textbf{Diff sequence:} Seven consecutive difference images
(shape $7 \times 1 \times 128 \times 128$) are each encoded
by a shared \textbf{DiffFrameCNN} (convolutional encoder,
64-dim output) and fed into a 2-layer \textbf{LSTM}
(hidden = 64). Both the final hidden state and the mean
of all hidden states are retained ($2 \times 64 = 128$-dim).

\textbf{Last raw frame:} The final frame of the window is
encoded by a separate \textbf{FrameCNN} to 32-dim, providing
instantaneous pool morphology context.

\textbf{Height patch (gated):} For the five profilometry
layers, a $33 \times 51$ pixel rectangle is extracted from
the height map $Z$ centered at the frame's aligned
$(x, y_\text{track})$ position, covering $\pm 0.197$\,mm
along scan and $\pm 0.200$\,mm cross-track — the full
melt pool footprint in both directions. The patch is
z-score normalized and downsampled to $16 \times 16$ before
encoding by \textbf{HeightPatchCNN} to 16-dim. The branch
output is multiplied by a binary gate $g$:
\begin{equation}
\mathbf{h}_\text{feat} = \text{HeightPatchCNN}(\text{patch}) \times g
\end{equation}
where $g = 1$ for L248--L252 and $g = 0$ otherwise.
This gating is essential: feeding zero patches without gating
produces spurious nonzero CNN outputs (due to bias terms and
BatchNorm) that mislead the classifier. The gate scalar $g$
is also concatenated explicitly into the fused vector so
the classifier can distinguish missing data from a genuinely
flat surface.

\textbf{Process parameters:} Normalized laser power $P$ and
scan speed $V$ are encoded by a linear branch to 16-dim.

The fused vector is:
\begin{equation}
\mathbf{z} = [\,\mathbf{h}_\text{LSTM}^{64},\;
\bar{\mathbf{h}}_\text{LSTM}^{64},\;
\mathbf{f}_\text{frame}^{32},\;
\mathbf{h}_\text{feat}^{16} \cdot g,\;
\mathbf{f}_{PV}^{16},\;
g^1\,] \in \mathbb{R}^{193}
\end{equation}
followed by a two-layer classifier (Linear(193, 64) $\to$
ReLU $\to$ Linear(64, 1)).

\begin{figure}[H]
\centering
\fbox{\rule{0pt}{5cm}\rule{0.9\linewidth}{0pt}}
\caption{Model architecture overview. Four input streams
are fused into a 193-dim vector before classification.
The height branch is gated to zero for layers without
profilometry data.}
\label{fig:arch}
\end{figure}

\subsection{Training and Evaluation}

\textbf{LOLO evaluation:} 14-fold leave-one-layer-out
cross-validation. Each fold trains on 12 layers, validates
on 1, and tests on the held-out layer. Validation layer
selection uses the layer most similar in $(P, V)$ space
to the test layer:
\begin{equation}
d(l, \text{test}) =
\left(\frac{P_l - P_\text{test}}{140}\right)^2 +
\left(\frac{V_l - V_\text{test}}{200}\right)^2
\end{equation}
This reduces checkpoint selection noise from mismatched
process conditions that would arise with random validation
layer selection.

\textbf{Training:} Adam optimizer, lr = $10^{-3}$,
weight decay = $10^{-3}$, batch size = 8, max 60 epochs
with early stopping (patience = 10) on validation AUC.
No data augmentation (tested but did not improve results
at this dataset size).

\textbf{Primary metric:} AUC-ROC, threshold-independent.
We report both mean per-layer AUC (averaging 14 individual
fold AUCs equally) and pooled AUC (computed on all
predictions combined), which answer different questions:
the former measures generalization across process conditions;
the latter measures overall detection capability weighted
by stratum size.

% ══════════════════════════════════════════════════════════════════
\section{Results}
% ══════════════════════════════════════════════════════════════════

\subsection{Per-Layer Performance}

\begin{table}[H]
\centering
\caption{Per-layer LOLO AUC. Camera-only baseline vs best
multi-modal model (gated height patch, smart val, SEQ=8).
[H] = layer receives real height patch at test time.}
\begin{tabular}{cccccc}
\toprule
Layer & P (W) & V (mm/s) & Camera only & Multi-modal & $\Delta$ \\
\midrule
L226 & 360 & 1800 & 0.640 & 0.707 & $+0.067$ \\
L227 & 380 & 1800 & 0.641 & 0.736 & $+0.095$ \\
L228 & 400 & 1800 & 0.658 & 0.782 & $+0.124$ \\
L229 & 420 & 1800 & 0.522 & 0.529 & $+0.007$ \\
L230 & 440 & 1800 & 0.594 & 0.576 & $-0.018$ \\
L231 & 460 & 1800 & 0.652 & 0.589 & $-0.063$ \\
L245 & 320 & 2000 & 0.659 & 0.654 & $-0.005$ \\
L246 & 340 & 2000 & 0.760 & 0.834 & $+0.075$ \\
L247 & 360 & 2000 & 0.672 & 0.601 & $-0.071$ \\
\midrule
L248 [H] & 380 & 2000 & 0.608 & 0.827 & $+0.219$ \\
L249 [H] & 400 & 2000 & 0.844 & 0.809 & $-0.036$ \\
L250 [H] & 420 & 2000 & 0.779 & 0.822 & $+0.043$ \\
L251 [H] & 440 & 2000 & 0.585 & 0.786 & $+0.201$ \\
L252 [H] & 460 & 2000 & 0.624 & 0.824 & $+0.201$ \\
\midrule
\textbf{Mean (all)} & & & 0.673 & \textbf{0.720} & $+0.047$ \\
\textbf{Mean [H]}   & & & 0.688 & \textbf{0.814} & $+0.126$ \\
\textbf{Mean no-H}  & & & 0.640 & 0.668 & $+0.023$ \\
\bottomrule
\end{tabular}
\end{table}

\begin{figure}[H]
\centering
\includegraphics[width=\linewidth]{lolo_enhanced_comparison.png}
\caption{Per-layer AUC comparison and delta. Layers with
profilometry data ([H]) show the largest improvements;
the non-profilometry layers show modest improvements from
smart validation selection and longer sequence length.}
\label{fig:lolo_comparison}
\end{figure}

All five profilometry layers improve, with improvements largest
on previously weak layers (L248: $+0.219$; L251, L252: $+0.201$)
and smallest on the already-strong L249 ($-0.036$, a ceiling
effect at baseline AUC = 0.844). The nine non-profilometry
layers show a small mean improvement (+0.023) driven by the
smart validation selection and longer sequence window rather
than the height patch.

\subsection{Per-Stratum Performance}

\begin{table}[H]
\centering
\caption{Per-stratum pooled AUC across all 14 LOLO folds.
Pooled AUC is computed on all predictions combined and differs
from mean per-layer AUC (camera-only: 0.673; multi-modal: 0.720)
which averages 14 individual fold AUCs equally.
$N$ = balling windows in stratum.}
\begin{tabular}{lrrrr}
\toprule
Stratum & $N$ & Camera only & Multi-modal & $\Delta$ \\
\midrule
Overall      & 1441 & 0.611 & 0.673 & $+0.062$ \\
\midrule
Type=Middle  &  492 & 0.624 & 0.750 & $+0.126$ \\
Type=Left    &  158 & 0.643 & 0.719 & $+0.076$ \\
Type=Right   &   93 & 0.804 & 0.760 & $-0.044$ \\
\midrule
Size=L       &   51 & 0.789 & 0.855 & $+0.066$ \\
Size=M       &  200 & 0.698 & 0.793 & $+0.095$ \\
Size=S       &  492 & 0.616 & 0.713 & $+0.097$ \\
\bottomrule
\end{tabular}
\end{table}

The largest stratum gains are for Middle-type (+0.126) and
Size=S (+0.097) beads — precisely the categories that the
camera cannot reliably detect alone. Middle beads produce no
lateral camera asymmetry; the height patch provides the only
available signal for them. Small beads ($< 0.10$\,mm) cause
sub-resolution camera perturbations; the profilometry detects
the physical bead regardless of camera resolution. Right-type
beads degrade slightly ($-0.044$): since they have lower
profilometry height than clean frames ($d = -0.378$), the height
patch adds a slightly misleading signal for this category —
physically consistent with the camera detecting gas-flow-driven
pool dynamics rather than bead height. Right-type beads remain
the most detectable category overall (AUC = 0.760).

% ══════════════════════════════════════════════════════════════════
\section{Analysis}
% ══════════════════════════════════════════════════════════════════

\subsection{Contribution of Each Model Component}

Table~\ref{tab:ablation} shows the incremental contribution of
each design decision by comparing progressively complete versions
of the model.

\begin{table}[H]
\centering
\caption{Ablation of model components. Each row adds one
design decision to the previous. Profil/no-H split by
profilometry availability at test time.}
\label{tab:ablation}
\begin{tabular}{lccccc}
\toprule
Configuration & Mean AUC & $\Delta$ & Profil $\Delta$ & No-H $\Delta$ \\
\midrule
Camera only, SEQ=6, random val   & 0.673 & ---    & ---    & --- \\
+ height patch (no gate)         & 0.686 & +0.013 & +0.120 & $-0.026$ \\
+ gate + smart val (SEQ=6)       & 0.702 & +0.029 & +0.100 & +0.010 \\
+ SEQ=8                          & \textbf{0.720} & \textbf{+0.047} & \textbf{+0.126} & \textbf{+0.023} \\
\bottomrule
\end{tabular}
\end{table}

The ungated height patch already improves profilometry layers
(+0.120) but slightly degrades non-profilometry layers ($-0.026$)
because zero-valued patches produce spurious nonzero CNN outputs
that confuse the classifier. Adding the explicit gate and smart
validation selection recovers this degradation and improves
non-profilometry layers (+0.010), confirming that both design
decisions independently contribute. Increasing SEQ from 6 to 8
frames adds a further +0.018 across the board, consistent with
capturing more complete balling events (75th percentile duration:
7 frames).

\subsection{Process Parameter Contribution}

To quantify the contribution of laser power $P$ and scan speed
$V$, we trained four variants — P+V, P-only, V-only, and no
process params — under otherwise identical conditions.

\begin{table}[H]
\centering
\caption{Process parameter ablation. Mean LOLO AUC and marginal
contributions. dP = AUC(P+V)$-$AUC(V-only);
dV = AUC(P+V)$-$AUC(P-only).}
\begin{tabular}{lcccc}
\toprule
Variant & Mean AUC & P adds over V & V adds over P \\
\midrule
P+V  & 0.677 & $-0.001$ & $-0.012$ \\
P-only & \textbf{0.689} & --- & --- \\
V-only & 0.677 & --- & --- \\
None   & 0.669 & --- & --- \\
\bottomrule
\end{tabular}
\end{table}

P-only outperforms P+V (0.689 vs 0.677), and adding $V$ to
$P$ marginally degrades performance ($-0.012$). $V$ takes
only two values in the dataset (1800 and 2000\,mm/s), making
it a binary group indicator. The model can exploit $V$ as a
shortcut — learning categorical detection rules per speed
group rather than generalizable pool dynamics — which hurts
on layers where the visual evidence conflicts with the speed
cue (L227: $-0.092$; L229: $-0.094$; L250: $-0.081$). $P$
varies continuously across six levels, making such shortcut
learning less likely.

Notably, for the strongest profilometry layers (L249, L250,
L252), removing all process parameters outperforms P+V (by
$+0.015$, $+0.098$, $+0.073$ respectively). The profilometry
height patch encodes the physical \emph{outcome} of the
process conditions — bead geometry — more directly than
the input conditions $P$ and $V$ themselves. When height
data is available, explicit process parameter conditioning
becomes redundant.

\begin{figure}[H]
\centering
\includegraphics[width=\linewidth]{pv_ablation.png}
\caption{Process parameter ablation. Left: per-layer AUC
for all four variants. Right: delta AUC relative to P+V —
positive values indicate P+V outperforms the ablated variant.
The inconsistent pattern of P+V vs P-only across layers
reflects the shortcut learning effect of the binary $V$
variable.}
\label{fig:pv_ablation}
\end{figure}

\subsection{Vision-Language Model Baseline}

As a zero-shot and few-shot baseline, we evaluated Claude
Sonnet on 5-frame sequence grids of difference images —
the same temporal context as the CNN-LSTM.

\begin{table}[H]
\centering
\caption{VLM vs CNN-LSTM AUC on five representative layers.}
\begin{tabular}{ccccccc}
\toprule
Layer & P & V & CNN-LSTM & Zero-shot & Few-shot \\
\midrule
L227 & 380 & 1800 & 0.711 & \textbf{0.700} & 0.453 \\
L247 & 360 & 2000 & 0.705 & 0.535 & 0.536 \\
L249 & 400 & 2000 & 0.842 & 0.610 & 0.561 \\
L250 & 420 & 2000 & 0.711 & 0.560 & 0.690 \\
L230 & 440 & 1800 & 0.569 & 0.512 & 0.512 \\
\midrule
Mean & & & \textbf{0.708} & 0.583 & 0.550 \\
\bottomrule
\end{tabular}
\end{table}

On the highest-signal layer (L227, Cohen's $d = +0.646$),
zero-shot VLM achieves AUC = 0.700 vs CNN-LSTM 0.711 — a
gap of only 0.011 — demonstrating that when balling is
visually obvious, a general-purpose foundation model with
no domain-specific training is competitive. On lower-signal
layers requiring subtle temporal pattern recognition (L249,
L250), the gap widens substantially (up to 0.232).

Counterintuitively, few-shot examples from other process
conditions hurt on average ($0.583 \to 0.550$). Melt pool
appearance is strongly process-condition-dependent: balling
at P = 380\,W, V = 1800\,mm/s looks sufficiently different
from balling at P = 400\,W, V = 2000\,mm/s that cross-condition
examples confuse rather than calibrate the model. This directly
motivates explicit process parameter conditioning in trained
models. With process-condition-matched few-shot examples (same
$P$ and $V$ group as the test layer), VLM performance would
likely improve substantially. More fundamentally, the VLM
is operating at a severe data disadvantage: the CNN-LSTM
is trained on thousands of labeled frames from matched
conditions, while the VLM receives at most a handful of
examples. At larger dataset scales (hundreds of labeled
layers), in-context learning with many matched examples
or fine-tuned VLMs may close this gap further.

% ══════════════════════════════════════════════════════════════════
\section{Discussion and Future Work}
% ══════════════════════════════════════════════════════════════════

The central finding of this work is that melt pool camera
and surface profilometry are genuinely complementary sensing
modalities for balling detection — they measure different
physical phenomena and their fusion substantially improves
detection of the categories that were previously undetectable
from camera alone. Middle-type beads (the majority of balling
events at 64\%) and small beads go from near-chance to
AUC $\approx 0.75$ and 0.71 respectively when height data
is incorporated.

The most immediate next step is obtaining profilometry for
the remaining nine layers (L226--L231, L245--L247). With full
profilometry coverage, several additional improvements become
testable: local anomaly normalization of the height patch
(subtracting local neighbourhood mean rather than global
track mean, making the patch invariant to substrate tilt)
improved profilometry layers by an additional +0.016 over
global normalization in preliminary experiments but was
limited by the nine zero-patch layers. With full coverage,
this representation is expected to outperform global
normalization across all folds.

The process parameter ablation suggests running the best
model with P-only conditioning (removing $V$) as a direct
next step — a one-parameter change that may further improve
the mean AUC given P-only already outperforms P+V in the
ablation (0.689 vs 0.677).

More broadly, this work demonstrates a proof of concept for
multi-modal in-situ/ex-situ fusion in LPBF monitoring. The
finding that profilometry height is nearly uncorrelated with
camera prediction scores ($r \approx 0$) confirms that
optical emission and geometric outcome are measuring
physically distinct aspects of the balling phenomenon,
and that both are needed for complete detection.

% ══════════════════════════════════════════════════════════════════
\section{Conclusion}
% ══════════════════════════════════════════════════════════════════

We presented a multi-modal CNN-LSTM for in-situ balling
detection in LPBF that fuses melt pool camera image sequences
with spatially aligned profilometry height patches. The model
achieves mean LOLO AUC of 0.720 (+0.047 over camera-only
baseline) and mean AUC of 0.814 on the five profilometry
layers (+0.126). Key contributions include: an explicit
gating mechanism for missing profilometry data, smart
validation layer selection by process condition proximity,
and an empirical demonstration that camera and profilometry
are complementary rather than redundant — the two sensors
are nearly uncorrelated and their fusion recovers detection
of bead categories that are fundamentally undetectable
from camera alone. The gas flow asymmetry finding provides
a physical explanation for the Left/Right detectability
difference and a concrete recommendation for future
experimental design: alternating scan direction between
layers eliminates the systematic atmospheric bias.

\end{document}
